\section*{1. Kết quả đạt được của Đề tài}
Qua quá trình nghiên cứu, phân tích, thiết kế và hiện thực hóa đề tài \textbf{“Hệ thống Giám sát và Bảo trì Cơ sở Hạ tầng Tích hợp Thực tế Tăng cường (AR-IMMS)”}, nhóm chúng em đã hoàn thành toàn diện các mục tiêu học thuật và thực tiễn đề ra, giải quyết thành công bài toán kết nối giữa dữ liệu vận hành số và không gian phần cứng vật lý:
\begin{itemize}[leftmargin=2em]
    \item \textbf{Xây dựng hoàn chỉnh mô hình quản trị hạ tầng số:} Thiết kế thành công cây phân cấp không gian đa tầng (\textit{Site $\rightarrow$ Room $\rightarrow$ Rack $\rightarrow$ Server $\rightarrow$ Workload/Container}) kết hợp mô hình Digital Twin, giúp số hóa toàn bộ sơ đồ không gian trung tâm dữ liệu.
    \item \textbf{Hiện thực hóa hệ thống thu thập và giám sát thời gian thực:} Xây dựng thành công hạ tầng Mini Data Center testbed gồm các máy chủ vật lý và tác nhân thu thập nhẹ (\textit{Collector Agent}), kết hợp với NestJS backend và WebSocket (Socket.IO) để truyền tải thông số telemetry liên tục, cho phép điều hướng phân cấp trực quan từ vĩ mô đến chi tiết thông số container.
    \item \textbf{Tích hợp thành công công nghệ Thực tế Tăng cường (AR):} Ứng dụng di động AR trên nền tảng React Native tích hợp mã định danh không gian (\textit{QR/ArUco Code}) cho phép kỹ sư hiện trường quét thiết bị vật lý để tức thời hiển thị lớp phủ thông tin trạng thái hoạt động, chỉ số tài nguyên và cảnh báo trực tiếp trên camera.
    \item \textbf{Hoàn thiện hệ sinh thái quản lý vòng đời và sự cố:} Xây dựng đầy đủ các phân hệ quản lý cảnh báo ngưỡng tự động, tiếp nhận và điều phối phiếu bảo trì (\textit{Tickets}), kiểm soát phân quyền RBAC cùng hệ thống nhật ký kiểm toán bất biến (\textit{Audit Trail}) phục vụ tuân thủ vận hành doanh nghiệp.
\end{itemize}

\section*{2. Đánh giá, Hạn chế và Hướng phát triển}
Mặc dù hệ thống AR-IMMS đã đáp ứng xuất sắc các tiêu chí về hiệu năng thời gian thực, tính sẵn sàng, bảo mật thông qua mTLS/JWT và khả năng mô phỏng trực quan, đồ án vẫn tồn tại một số giới hạn nhất định trong phạm vi nghiên cứu:
\begin{itemize}[leftmargin=2em]
    \item \textbf{Phạm vi môi trường kiểm thử:} Do giới hạn về nguồn lực phần cứng, hệ thống hiện mới chỉ được đánh giá trên mô hình thử nghiệm cục bộ (\textit{Mini Data Center testbed}) với quy mô node và container giả lập, chưa thể hiện thực hóa trên quy mô hàng nghìn máy chủ thực tế của một trung tâm dữ liệu lớn.
    \item \textbf{Tính năng tự động hóa nâng cao:} Các thuật toán học máy dự báo sớm sự cố phần cứng hoặc tự động gán nhãn thông minh dựa trên lịch sử telemetry chưa được tích hợp chuyên sâu vào tầng xử lý cảnh báo.
\end{itemize}

\textbf{Hướng phát triển trong tương lai:} Nhóm sẽ tiếp tục mở rộng và nâng cấp hệ thống theo các định hướng sau:
\begin{itemize}[leftmargin=2em]
    \item Tích hợp cơ chế phân vùng dữ liệu nâng cao (\textit{Table Partitioning}) và tối ưu hóa lưu trữ chuỗi thời gian (\textit{Time-series data optimization}) nhằm xử lý khối lượng lớn metric từ hàng nghìn thiết bị.
    \item Phát triển chế độ đồng bộ ngoại tuyến (\textit{Offline-first sync}) cho ứng dụng di động AR nhằm hỗ trợ kỹ thuật viên thao tác ổn định trong các khu vực hạ tầng hầm máy chủ bị mất kết nối mạng.
    \item Nghiên cứu tích hợp các mô hình Trí tuệ nhân tạo (AI) để tự động phân tích nguyên nhân gốc rễ của sự cố (\textit{Root Cause Analysis}) và đề xuất hướng xử lý trực tiếp lên không gian AR.
\end{itemize}
